\section{방법론}
\label{sec:method}

\subsection{문제 정식화}
\label{sec:method:formulation}

% P1: 예측 시점, 특징 절단 시점, 위험집합, 라벨
경기 $m$에서 단계 $k$의 종료 시각을 $b_{m,k}$, 특징 절단 시각을 $a_{m,k}=b_{m,k}-\Delta$($\Delta=10$초)라 하자.
팀 $i$의 탈락 시각을 $T_{m,i}$, 경기 종료 시각을 $C_m$이라 하면, 관측되는 시각은 $\widetilde T_{m,i}=\min(T_{m,i},C_m)$이고 탈락 여부는 $\delta_{m,i}=\mathbb{1}[T_{m,i}\le C_m]$이다. 우승 팀은 $\delta_{m,i}=0$이다.
예측 시점은 단계 경계 $b_{m,k}$이며, 위험집합은 이 시점에 생존한 팀들의 집합 $\mathcal{R}_{m,k}=\{i:\widetilde T_{m,i}>b_{m,k}\}$이다.
위험집합의 각 팀에 대해 라벨은 $Y_{m,i,k}=\mathbb{1}[b_{m,k}<T_{m,i}\le b_{m,k+1}]$이며, 예측 대상은 다음 단계 탈락 위험
\begin{equation}
\lambda_{m,i,k}=\Pr\!\left(T_{m,i}\le b_{m,k+1} \mid T_{m,i}>b_{m,k},\,\mathcal{H}_{m}(a_{m,k})\right)
\label{eq:hazard}
\end{equation}
이다. 여기서 $\mathcal{H}_{m}(t)$는 시각 $t$까지의 경기 기록이다.
다음 단계가 경기의 마지막 단계이면 $b_{m,k+1}=C_m$으로 둔다.
위험집합의 구성에는 예측 시점의 생존 여부를 사용하지만, 모델의 입력 특징, 생존 인원, 팀 간 관계, attention 마스크는 특징 절단 시각 이전의 정보로만 구성한다.
이 정의는 각 단계 경계를 landmark로, 다음 단계를 예측 구간으로 하는 이산시간 hazard에 해당한다.
본 연구는 hazard를 로짓 척도에서 단계별 기저위험과 팀 상태의 함수로 모형화한다.
\begin{equation}
\mathrm{logit}\,\lambda_{m,i,k}=\alpha_k+f(\mathbf{x}_{m,i,k},\mathcal{G}_{m,k})
\label{eq:logit}
\end{equation}
여기서 $\mathbf{x}_{m,i,k}$는 팀의 서식지 맥락, $\mathcal{G}_{m,k}$는 같은 단계의 다른 팀들과의 관계이다.

% P2: 상대위험 학습 목적과의 차이
본 연구에서 비교하는 상대위험 학습 목적은 같은 위험집합 안에서 팀들의 점수 차이를 학습한다.
표준 Cox 부분우도의 각 항은 사건 팀의 점수가 위험집합에서 차지하는 비율 $\exp(s_i)/\sum_{j}\exp(s_j)$이며, KRAFTON의 weighted Cox는 이를 가중한 변형이다.
이러한 목적은 같은 시점의 모든 점수에 공통으로 더해지는 이동에 불변이므로, 그 자체로는 식~\eqref{eq:logit}의 $\alpha_k$와 같은 단계별 공통 위험 수준을 식별하지 못한다.
기저위험을 따로 추정하여 결합하면 구간별 확률을 산출할 수 있으므로, 본 연구는 상대위험 점수에 단계별 기저위험을 결합한 변형도 함께 비교한다(\ref{sec:method:decomposition}절).

\subsection{서식지 맥락 표현}
\label{sec:method:representation}

% P3: 구성개념 네 가지
서식지 축소 관점에서 각 팀의 상태를 네 가지 구성개념으로 정의한다.
첫째, 구역 경계에 대한 노출은 팀 중심과 안전구역 경계 사이의 부호 있는 거리로 측정하며, 구역 밖에 있는 깊이와 구역 안쪽으로 들어온 깊이로 나누어 사용한다.
둘째, 이동 부담은 안전구역에 도달하는 데 필요한 속도와 팀이 실제로 이동한 속도의 차이, 즉 이동 부족분으로 측정한다~\cite{Nathan08}.
필요 속도는 특징 절단 시각 직전의 팀 중심에서 그 시점에 공개된 안전구역 경계까지의 거리(구역 안이면 0)를 단계 길이로 나눈 값이고, 이동 속도는 단계 시작부터 특징 절단 시각까지의 이동 경로 길이를 같은 단계 길이로 나눈 값이다.
단계 길이에는 경기별 관측값 대신 학습 집합에서 코호트, 경기 유형, 단계별로 구한 중앙값을 사용하여, 특징 절단 시각 이후의 정보가 들어가지 않게 하였다.
셋째, 주변 경쟁 밀도는 팀 주변 200\,m 안에 있는 다른 팀의 수이다.
넷째, 남은 팀원 수는 특징 절단 시각에 생존한 것으로 관측된 팀원의 수이다.
학습 집합에서 필요 속도는 구역 밖 깊이와 순위 상관이 0.82로 높아 정보가 중복될 가능성이 있으므로, 실제 이동을 함께 반영하는 이동 부족분(상관 0.38)을 사용하였다.
이 선택은 결과를 확인한 뒤 이루어진 탐색적 특징 설계이다.

% P4: 전체 특징 표현과 관계
예측 모델의 입력은 이 구성개념을 포함하는 150차원의 팀--단계 특징이다.
특징은 경기 맥락, 팀 상태, 안전구역, 이동, 노출, 자원, 공간, 관계, 서식지의 아홉 계열로 구성되며, 여기에 팀 중심 좌표와 팀의 AI 계정 비율을 더한다.
같은 단계의 팀들 사이에는 세 유형의 관계를 정의한다.
근접 관계는 단계 시작부터 특징 절단 시각까지 두 팀의 평균 거리 $d$에 대해 가중치 $\exp(-d/200)$를 갖는 대칭 관계이다.
교전 관계는 같은 구간에서 한 팀이 다른 팀에 가한 피해, 기절, 처치를 결합한 강도 $\log(1+\text{강도})$를 가중치로 갖는 방향 관계이며, 교전이 있었던 팀 쌍에서는 근접 관계를 교전 관계로 대체한다.
공동 압박 관계는 두 팀이 모두 구역 밖에 있거나 경계 안쪽 150\,m 이내에 있으면서 서로 200\,m 이내에 있을 때의 대칭 관계이다.

\subsection{RQ1: 생태 hazard 모델}
\label{sec:method:rq1}

% P5: 해석 가능한 모델과 일관성 평가
RQ1에 답하기 위해 \ref{sec:method:representation}절의 구성개념만으로 식~\eqref{eq:logit}을 로지스틱 회귀로 추정한다.
\begin{equation}
\mathrm{logit}\,\lambda_{m,i,k}=\alpha_k+\sum_{c}\beta_c z_{c,m,i,k}
\label{eq:ehm}
\end{equation}
여기서 $z_c$는 학습 집합에서 표준화한 구성개념(구역 밖 깊이, 구역 안쪽 깊이, 이동 부족분, 주변 경쟁 밀도, 남은 팀원 수)이며, 전체 생존 팀 수의 로그를 통제 변수로 포함한다.
솔로 코호트에서는 남은 팀원 수가 상수이므로 제외한다.
모델은 코호트마다 학습 집합에서 추정하며, 표준오차는 경기 단위로 군집화하여 계산한다.
코호트별 계수는 DerSimonian--Laird 확률효과 모형으로 통합하고, 통합 추정치와 95\% 신뢰구간, 코호트 간 이질성($\tau^2$, $I^2$), 그리고 계수의 부호가 통합 추정치와 같은 코호트의 수를 보고한다.
단계에 따른 변화를 보기 위해 각 구성개념과 단계 구간(초반 1--2, 중반 3--5, 후반 6--7)의 상호작용을 추가한 모델을 추정하고, 구간별 기울기와 구간 간 기울기가 같은지에 대한 검정 결과를 보고한다.

\subsection{RQ2: 비교 설계}
\label{sec:method:rq2}

% P6: 본 연구의 설계
본 연구의 설계는 식~\eqref{eq:logit}의 단계 조건부 hazard를 다음 단계 탈락에 대한 교차엔트로피로 학습하고, \ref{sec:method:representation}절의 서식지 맥락 표현을 입력으로 사용한다.
설계가 특정 학습기에 의존하지 않음을 보이기 위해 두 학습기로 구현한다.
첫째는 단계 경계의 팀 특징에 대한 gradient boosting 분류기(HGB; 트리 200개, 학습률 0.06)로, 단계 번호를 입력 특징으로 포함하며 식~\eqref{eq:logit}의 $\alpha_k$를 별도의 항으로 두지는 않는다.
둘째는 참조 모델인 PTR-Former로, 팀--단계 토큰에 단계 임베딩을 더한 뒤 같은 팀의 단계 이력에 대한 인과적 attention(단계 축)과 같은 단계의 생존 팀들에 대한 attention(팀 축)을 번갈아 적용하며, 팀 축 attention의 로짓에는 관계 유형별로 학습되는 편향을 더한다.
PTR-Former는 다음 단계 hazard의 교차엔트로피에 최종 순위에 대한 쌍별 순위 손실을 0.1의 가중치로 더하여 학습한다\footnote{차원 128, attention 헤드 4개, 블록 2개, dropout 0.1을 사용하였다. AdamW(학습률 $10^{-3}$, 가중치 감쇠 $10^{-4}$)와 코사인 학습률 감소로 최대 12 epoch 동안 경기 8개씩 묶어 학습하였으며, 검증 집합의 로그 손실이 3 epoch 동안 개선되지 않으면 학습을 멈추고 가장 좋았던 시점의 가중치를 사용하였다.}.

% P7: 기존 설계의 적용
두 기존 설계는 원 설계의 입력과 학습 목적을 유지하면서 본 연구의 표본에 맞게 적용하였으며, 적용 내용은 표~\ref{tab:adapt}에 정리하였다.
Sagong 등의 설계~\cite{Sagong25a}는 플레이어 단위의 활동 지표로 구성되므로 2인·4인 코호트에서는 팀 단위로 집계하였고, 본 연구의 정규화 자료에 포함되지 않은 공격 시도와 어시스트 기록은 가장 가까운 사건으로 대체하였다.
KRAFTON의 예측기~\cite{KRAFTON}는 공개 구현(commit 99ed9d1)을 기반으로, 입력 지표를 본 연구의 텔레메트리로 재구성하고 공개 경기와 평가 과제에 맞춘 적용판으로 구성하였다.
모든 설계는 같은 경기 분할을 사용하며, 평가는 같은 팀--단계 위험집합과 다음 단계 탈락 라벨에서 수행한다.
RQ2에서는 각 적용판의 원래 학습 표집과 감독 목적을 유지하고, RQ3에서는 학습 표본과 입력을 추가로 통제한다.

% P7b: 적용판의 충실도 점검
적용판이 본 연구의 과제에 불리하게 구성되지 않았는지 확인하기 위해, 각 적용판을 원래 과제의 지표로도 평가한다.
KRAFTON 적용판은 원 구현이 정의한 우승 팀 정확도와 탈락 순서의 C-index로, Sagong 적용판은 원 논문과 같은 개인 단위의 생존 분류로 평가한다.
본 연구의 설계도 같은 지표로 평가하여, 과제가 바뀔 때 설계 간 우열이 어떻게 달라지는지 확인한다.

\begin{table}[t]
\centering
\caption{기존 설계의 적용}
\label{tab:adapt}
\small
\begin{tabular}{p{1.3cm}p{3.1cm}p{3.1cm}}
\hline
항목 & Sagong 적용판 & KRAFTON 적용판 \\
\hline
원 설계 & 플레이어 단위, 다음 반 단계 생존 분류 & 팀 단위, 사망 시각 Cox 학습 후 우승확률로 변환 \\
입력 & 단계별 활동 지표 8개, FightScore, ActiveScore, 플레이 유형, 유형별 전이 생존율 (12개) & 팀 지표 55개(원 정의와 같음 32, 근사 22, 재구성하지 못해 상수 0으로 대체 1), 팀원 위치, 구역·맵 임베딩 \\
학습기 & XGBoost, Random Forest, 로지스틱 회귀 & 8층 Transformer, weighted Cox \\
학습·평가 시점 & 단계 경계 / 단계 경계 & 원 설계의 시점 표집 / 단계 경계 \\
적용 결정 & 공격은 적에게 가한 피해 사건 수로, 어시스트는 다른 플레이어의 처치 직전 60초 안의 피해로 대체; 2·4인 팀은 횟수 합, 밀집도 최댓값, 이동 거리 평균으로 집계; 반 단계 대신 단계를 집계 구간으로 사용 & 팀 수 상한 제거, 탈락 팀의 attention 차단, 평가 코드의 점수 방향 해석 수정(99ed9d1), 검증 집합에서 Platt 보정(같은 검증 집합에서 단계별 기저위험을 추정한 +H0 변형 병행) \\
설정 & 라이브러리 기본값(원 논문 미기재) & 원 설계의 설정 \\
\hline
\end{tabular}
\end{table}

\subsection{RQ3: 단일 요인 분해}
\label{sec:method:decomposition}

% P8: 분해 사다리
RQ3에 답하기 위해 다른 조건을 고정한 채 한 번에 하나의 설계 요인만 바꾸는 비교를 구성한다(표~\ref{tab:ladder}).
입력 표현의 효과는 학습기를 HGB로 고정하고 Sagong 등의 입력 또는 재구성한 KRAFTON 입력을 본 연구의 입력으로 바꾸어 측정한다.
Sagong 등의 입력과의 비교는 서식지 맥락을 먼저 추가한 뒤 나머지 입력을 바꾸는 두 단계로 나누어, 차이 가운데 서식지 맥락이 차지하는 부분을 구분한다.
학습 목적과 단계별 기저위험의 효과는 KRAFTON의 Transformer 구조에 본 연구의 입력을 넣고 weighted Cox 목적으로 학습한 모델(KR-Cox)에서 출발한다.
KRAFTON 원본에서 KR-Cox로 넘어가는 비교는 입력과 학습 표본이 함께 바뀌는 묶음 비교이므로, 분해 단계와 구분하여 보고한다.
KR-Cox의 학습 목적을 다음 단계 탈락의 교차엔트로피로 바꾼 모델을 KR-Hazard, KR-Hazard에 단계별 로짓 절편 $\alpha_k$를 추가한 모델을 KR-Hazard+$\alpha$라 한다.
KR-Cox와 KR-Hazard는 구조, 입력, 학습 경기와 절단 시점, 초기 가중치, 갱신 횟수가 같고 학습 목적만 다르다.
다만 KR-Cox는 특징 절단 시각에 생존한 팀들의 이후 탈락 순서를, KR-Hazard는 예측 시점의 위험집합에서 다음 단계의 탈락 여부를 감독 정보로 사용하므로, 이 비교는 원래의 상대위험 학습 목적과 과제에 맞춘 다음 단계 위험 학습 목적의 차이를 나타낸다.
입력에는 단계 번호가 포함되어 있으므로, KR-Hazard+$\alpha$와의 비교는 단계 정보를 포함한 입력을 유지한 상태에서 명시적인 단계별 로짓 절편을 추가하는 효과를 나타낸다.
Cox 점수 자체에는 단계별 위험 수준이 담기지 않으므로, KR-Cox의 점수에 학습 집합에서 추정한 단계별 기저위험을 결합한 KR-Cox+H0도 함께 평가한다~\cite{PrenticeGloeckler78}.
따라서 학습 목적의 비교는 KR-Cox에 기저위험을 결합하는 단계와, 기저위험을 갖춘 KR-Cox+H0를 KR-Hazard로 바꾸는 단계로 나누어 측정한다.
뒤의 단계는 기저위험을 갖춘 Cox 모델과 다음 단계 탈락확률을 직접 학습한 모델의 차이를 보여 준다.
아키텍처와 구조의 비교에서는 학습 목적을 맞추기 위해 순위 보조 손실을 사용하지 않는다.
아키텍처의 효과는 같은 학습 목적과 입력에서 KRAFTON 구조를 PTR-Former의 팀 축만 남긴 모델(PTR-Team)로 바꾸어 측정하고, 구조의 효과는 PTR-Team에 단계 축과 관계 편향을 추가한 PTR-Former와 비교하여 측정한다.
순위 보조 손실의 효과는 같은 구조에서 이를 추가한 모델과 비교하여 별도로 보고한다.
구조 요인은 단계 축과 팀 축을 2$\times$2로 켜고 끄는 요인 실험으로 평가하고, 관계의 대응을 경기 안에서 무작위로 섞어 다시 학습한 대조 모델로 관계 편향의 기여를 확인한다.
또한 안전구역의 원시 기하(현재·다음 구역의 중심과 반경)만을 갖는 구역 객체를 토큰으로 추가하여, 수작업 서식지 특징을 뺀 입력에서 그 정보를 대신 학습하는지 확인한다.
마지막으로 HGB의 입력에 재구성한 KRAFTON 지표를 추가하여 두 표현의 보완성을 확인한다.

\begin{table}[t]
\centering
\caption{단일 요인 분해의 단계}
\label{tab:ladder}
\small
\begin{tabular}{p{0.5cm}p{2.9cm}p{1.9cm}p{1.9cm}}
\hline
단계 & 비교 & 바꾸는 요인 & 고정하는 조건 \\
\hline
(a1) & HGB: Sagong 입력 $\rightarrow$ +서식지 맥락 & 서식지 맥락 & 학습기, 학습 목적 \\
(a2) & HGB: Sagong 입력+서식지 맥락 $\rightarrow$ 본 연구 입력 & 나머지 입력 & 학습기, 학습 목적 \\
(b) & HGB: KRAFTON 입력 $\rightarrow$ 본 연구 입력 & 입력 표현 & 학습기, 학습 목적 \\
(c1) & KR-Cox $\rightarrow$ KR-Cox+H0 & 단계별 기저위험 결합 & 구조, 입력, 표본, 학습 목적 \\
(c2) & KR-Cox+H0 $\rightarrow$ KR-Hazard & 학습 목적 & 구조, 입력, 표본 \\
(d) & KR-Hazard $\rightarrow$ KR-Hazard+$\alpha$ & 명시적 단계 절편 & 구조, 입력, 학습 목적 \\
(e) & KR-Hazard+$\alpha$ $\rightarrow$ PTR-Team & 아키텍처 & 학습 목적(순위 손실 없음), 입력 \\
(f) & PTR-Team $\rightarrow$ PTR-Former & 단계 축, 관계 편향 & 학습 목적(순위 손실 없음), 입력, 팀 축 \\
(g) & HGB: 본 연구 입력 $\rightarrow$ +KRAFTON 입력 & 입력 확장 & 학습기, 학습 목적 \\
\hline
\end{tabular}
\end{table}

\subsection{평가 프로토콜}
\label{sec:method:evaluation}

% P9: 지표
주 지표는 시험 집합의 위험집합 전체에 대한 AUC이다.
보완 지표로 두 가지를 함께 보고한다.
단계 안 AUC는 같은 단계 번호의 표본을 경기에 걸쳐 모아 단계별로 계산한 AUC의 평균으로, 단계마다 공통으로 더해지는 위험 수준의 차이를 제외한 판별력을 나타내며, 시험 표본이 30개 이상이고 두 라벨이 모두 있는 단계만으로 계산한다.
경기--단계 내부 AUC는 같은 경기의 같은 단계 위험집합 안에서 팀을 구분하는 판별력으로, 위험집합별 AUC를 양성--음성 쌍의 수로 가중하여 통합하고 라벨이 하나뿐인 위험집합은 제외한다.
로그 손실과 Brier 점수는 확률 예측의 전반적인 정확도를 나타낸다.
단계별 보정 오차는 단계마다 평균 예측확률과 관측 탈락률의 차이의 절댓값을 구해 표본 수로 가중 평균한 값으로, 설계가 단계마다 달라지는 위험 수준을 표현하는지를 보여 준다.
상대 순위에 대해서는 최종 순위와의 Spearman 상관과 상위 5위 판별 AUC를 사용한다.
상대 순위 지표에는 각 모델의 위험 점수를 사용한다. 순위 보조 손실로 학습한 PTR-Former는 순위 출력을, 나머지 모델은 다음 단계 탈락확률 또는 Cox 점수를 사용한다.

% P10: 통계적 추론과 분석의 구분
두 설계의 차이는 코호트마다 경기 단위로 2,000회 재표집한 쌍대 bootstrap으로 신뢰구간을 구하며, 이 구간은 학습된 모델을 고정한 상태에서 시험 표본이 달라질 때의 변동을 나타낸다.
12개 코호트의 결과는 DerSimonian--Laird 확률효과 모형으로 통합하고, REML과 Hartung--Knapp 방식의 결과를 민감도 분석으로 함께 확인한다.
의미 있는 최소 효과 크기는 AUC 0.005로 둔다.
통합 추정치의 95\% 신뢰구간 하한이 0보다 크고 점추정치가 0.005 이상이면 개선, 95\% 신뢰구간 상한이 0보다 작고 점추정치가 $-0.005$ 이하이면 저하, 90\% 신뢰구간이 $\pm0.005$ 안에 있으면 실질적 동등, 그 밖에는 불확정으로 분류한다.
개선 분류는 양의 차이에 대한 통계적 근거와 점추정치의 크기를 함께 요구하는 기준이며, 실제 효과가 0.005 이상임을 입증하는 검정은 아니다.
구조 요인의 다섯 비교는 실험 전에 타임스탬프가 남는 커밋으로 분석 계획을 고정한 사전 지정 분석이며, 다섯 개의 통합 검정에 Holm 방법을 적용한다\footnote{분석 계획은 2026년 10월 5일에 고정하였다. 다섯 비교는 단계 축과 팀 축의 2$\times$2 요인 실험에서 다른 축의 두 수준에 대해 평균한 단계 축의 주효과(AUC), 같은 방식의 팀 축 주효과(AUC), 팀 축 주효과의 Spearman 상관 차이, 관계의 대응을 경기 안에서 섞은 대조 모델 대비 관계 편향의 효과(AUC), 서식지 특징을 뺀 입력에서 구역 객체 토큰의 효과(AUC)이다. 계획을 고정한 뒤 데이터 처리의 누수를 발견하여 고쳤으며, 모든 결과는 고친 자료에서 다시 얻었다. 이 수정은 이전 결과를 확인한 뒤 결정하였다.}.
% 익명 심사본에서는 커밋 해시를 생략하고, 최종본에 추가한다(f59bac9).
생태 hazard 모델, 기존 설계와의 비교, 적용판의 원래 과제 평가, 단일 요인 분해, 단계 안 AUC, 경기--단계 내부 AUC, 학습 횟수 민감도, 학습 곡선은 이후에 추가한 탐색적 분석으로 구분하여 보고한다.
신경망 모델은 모든 코호트에서 seed 42로 학습하고, 가장 큰 두 코호트에서 seed 13과 2026을 추가하여 학습 초기화에 따른 변동을 확인한다.
seed 변동은 두 코호트에서만 확인하였으므로, 작은 코호트의 결과는 이 한계 안에서 해석한다.
학습 횟수의 상한이 신경망 모델의 성능을 제한하는지 확인하기 위해, 같은 두 코호트에서 상한을 24 epoch로 늘려 다시 학습한다.
데이터 규모의 영향은 같은 두 코호트에서 학습 경기 수를 300, 600, 1,200경기와 전체로 바꾼 학습 곡선으로 확인하며, epoch 수를 고정한 경우와 갱신 횟수를 맞춘 경우를 모두 구한다.
